\chapter{ThermoDynamics of Chemical Bonding}\label{ch:bonding}

\small
\begin{quote}
The bond is not a state. It is the trace, on the equilibrium surface, of a
process that took the system from one RealQM minimum to another and dissipated
the difference into the surrounding bath.
\end{quote}
\begin{quote}
\itshape The really fundamental thing is the interaction of electrons and
nuclei, which is governed by Coulomb's law. (Schr\"odinger, 1926)
\end{quote}
\normalsize

\section{The Bond as State vs the Bond as Event}

The classical thermodynamic description of a chemical bond consists of a few
state-function numbers --- a bond enthalpy $\Delta H_{\mathrm{bond}}$, a bond
entropy $\Delta S_{\mathrm{bond}}$, a free energy $\Delta G = \Delta H -
T\Delta S$ and the equilibrium constant $K = \exp(-\Delta G/RT)$ that decides
whether the bond exists at temperature $T$. The framework is silent on the
\emph{process} by which the bond forms. It encodes the released energy as a
state-function difference and the multiplicity loss as a state-function
difference, with the actual bonding event reduced to a discontinuous jump on
the equilibrium surface. There is no mechanism in it for how the energy is
released, or where it goes, or why an entropy decreases when two free atoms
become one bound molecule.

The view taken in this volume, and consistent with the with-Dynamics framing
of Chapter~\ref{ch:withwithout}, is different. The bond is not a state but
an \emph{event}: a process that takes the system from one RealQM minimum
(two separated atoms with their own electron territories) to another RealQM
minimum (a bonded molecule with overlapping electron territories between the
two nuclei), with the energy gap between the two minima dissipated into the
bath through the same molecular-scale Joule pattern that organises the rest
of this book. The state-function bookkeeping of classical thermodynamics
records the result; the dynamics records the process.

\section{The Atomic-Scale Energy Gap (RealQM Input)}

The energy gap between the two configurations is supplied by Volume~VI of
this programme, \emph{Real Quantum Mechanics} (RealQM)~\cite{RealQM2026}.
In the RealQM formulation, an $N$-electron system is described by a sum of
one-electron wave functions on non-overlapping spatial domains separated by
free (Bernoulli) boundaries, with ground states minimising the ordinary
Coulomb energy --- no configuration space, no antisymmetry postulate.

For the two-atom problem, two minima of the Coulomb energy functional
typically exist:

\begin{itemize}
\item A \emph{separated} minimum: the two atoms sit at large internuclear
distance $R \to \infty$ with their electron territories disjoint and
spherically symmetric about each nucleus. The total energy is
$E_{\mathrm{sep}} = E_A + E_B$, the sum of the isolated-atom energies.

\item A \emph{bonded} minimum: the two atoms sit at the equilibrium bond
distance $R_{\mathrm{eq}}$ with their electron territories rearranged so as
to overlap in the internuclear region, screening the proton--proton
repulsion and pulling both nuclei into the shared electron density. The
total energy is $E_{\mathrm{bond}}$, below $E_{\mathrm{sep}}$ by the
\emph{bond dissociation energy}
\begin{equation}
  D_e \;=\; E_{\mathrm{sep}} - E_{\mathrm{bond}} \;>\; 0.
  \label{eq:De}
\end{equation}
\end{itemize}

For $\mathrm{H}_2$ the RealQM calculation reproduced in Vol~VI gives $D_e
\approx 0.17\,\mathrm{Ha}$ at $R_{\mathrm{eq}} \approx 1.4\,\mathrm{au}$, in
agreement with the experimental value $4.75\,\mathrm{eV}$. The gap is a
deterministic property of the Coulomb energy functional --- it requires no
statistics, no probability interpretation, no microstate counting. Whatever
``bond energy'' the classical thermodynamic tables list for H--H, RealQM
computes from first principles.

\section{The Joule Pattern at Molecular Scale}

The bonding event, viewed as a dynamical process, follows the same
three-stage Joule pattern that organises Chapters \ref{ch:joule},
\ref{ch:refrig} and \ref{ch:piston}:

\begin{enumerate}
\item \textbf{Internal energy reservoir.} The two separated atoms together
hold an electronic configuration energy $E_{\mathrm{sep}}$. This is the
RealQM analogue of the high-pressure gas in Joule's left chamber: a stored
internal energy waiting to be released.

\item \textbf{Conversion to bulk kinetic energy.} As the atoms approach and
the electron territories rearrange, the energy gap $D_e$ is released by the
electronic system into the \emph{nuclear} degrees of freedom: first as
relative translational kinetic energy of the two approaching nuclei, then,
once they fall into the bonded minimum, as oscillatory kinetic energy of
the newly formed vibrational mode of the diatomic molecule. The energy is
not yet heat --- it is coherent kinetic motion of a single oscillator. This
is the analogue of the bulk jet through the channel in Joule's experiment:
internal energy converted into directed bulk flow.

\item \textbf{Dissipation into the bath.} The vibrational mode of the
newly formed molecule does not stay coherent. It couples to surrounding
molecules through collisions, transferring quanta of vibrational energy
into translational kinetic energy of the bath in a sequence of irreversible
inelastic collisions. After a few mean-free-path times, the coherent
vibration has been thermalised: the originally localised energy $D_e$ has
become a tiny incremental heating of the surrounding gas. This is the
analogue of the turbulent dissipation in the receiving chamber of Joule's
experiment: directed bulk kinetic energy degraded into incoherent thermal
motion.
\end{enumerate}

The result is the same as the result reported by classical thermodynamics
--- a heat release of magnitude $D_e$ per bond formed --- but the mechanism
is now explicit. The energy was not magically transferred from one column
of a table to another; it was released by the electronic system, channelled
into a single coherent mode of nuclear motion, and then dissipated into the
bath by collisional coupling. The Joule pattern --- internal energy
$\to$ bulk kinetic energy $\to$ irreversible dissipation $\to$ heat ---
applies at molecular scale exactly as it applies at macroscopic scale.

\section{Macroscopic Heat Balance and the Reactive RNS Apparatus}

For a gas in which the bonding event is happening many times per unit
volume per unit time, the bookkeeping is the business of the reactive RNS
apparatus developed in Chapter~\ref{ch:reactive} (Chemistry: RNS with
Chemical Reactions). Each elementary reaction $\mathrm{A} + \mathrm{B} \to
\mathrm{AB}$ contributes:

\begin{itemize}
\item A source term in the species equation for $\mathrm{AB}$ and matching
sink terms for $\mathrm{A}$ and $\mathrm{B}$, with rate set by collision
frequency and the Arrhenius factor for any activation barrier.

\item A heat source in the internal-energy equation equal to $D_e$ per
reaction event, contributing to the local temperature rise.

\item A momentum redistribution (typically small for a near-equilibrium
reaction) reflecting the change in mean molecular weight of the gas.
\end{itemize}

The viscous dissipation $D = \int\!\int\mu_{\mathrm{mom}}|\nabla u|^2\,dx\,dt$
of the bulk flow is now joined by a \emph{reactive dissipation} from the
local relaxation of species gradients --- the chemical analogue of viscous
shear. Together they form the irreversibility of a reacting gas: the 2nd
Law in this setting is the statement that both contributions are positive
and that their sum equals the rate at which the chemical free energy stored
in unreacted species is degraded to heat.

\section{The Mysteries of Bonding Entropy Dissolved}

Several classical puzzles about bonding now have natural resolutions:

\begin{itemize}
\item \emph{Why is $\Delta S_{\mathrm{bond}}$ negative?} In the classical
framework this is mysterious: the bonded molecule has fewer ``microstates''
than two free atoms, so $\Delta S$ must be negative, but the framework
gives no reason \emph{why} the bonded molecule is the natural endpoint
despite the entropic penalty. In the with-Dynamics framework the question
dissolves: bond formation is the dissipative dynamics relaxing toward the
deeper RealQM minimum, with the released energy $D_e$ dissipated into the
bath. There is no microstate counting. The bond exists because the
$D_e$-deep minimum exists, and the dissipative dynamics finds it.

\item \emph{Why is the equilibrium constant $K = \exp(-\Delta G/RT)$?} In
the with-Dynamics framework this is the steady-state condition for the
reactive RNS: at temperature $T$, the forward reaction rate (driven by
$D_e$) and the backward rate (driven by thermal kinetic energy of bonded
pairs sufficient to climb out of the bonded minimum) balance when the
ratio of populations is set by the Boltzmann factor evaluated at the
bonded-minimum depth. No postulate is needed; the steady state of the
reactive flow gives the same answer that the classical state-function
bookkeeping was forced to postulate.

\item \emph{Where does the bond energy ``go''?} Into the bath, by the
three-stage Joule pattern: into vibrational kinetic energy of the newly
formed molecule, and from there by inelastic collisions into translational
kinetic energy of surrounding molecules. The same dissipation $D > 0$
that organises the rest of this book carries the bonding event from RealQM
to the macroscopic temperature.
\end{itemize}

The chemical bond, the most basic non-trivial chemical event, is therefore
the cleanest single illustration of why ThermoDynamics needs Dynamics. The
classical state-function description captures the bookkeeping; the
with-Dynamics description captures the bookkeeping \emph{and} the
mechanism, with RealQM supplying the energy gap and the Joule pattern at
molecular scale supplying the channel through which the gap is dissipated.
